\subsection{Affine guardedness and covariance control}
\label{sec:leace}

Our subsequent results are formulated as minimal-disturbance affine transformations under constraints on the cross-covariance between representations and concept indicators. 
We use the guardedness framework of \citet{ravfogel-etal-2023-linear,belrose2025leace} to justify why cross-covariance is a meaningful object for formalizing linear concept information in representations. 
In particular, for the task of concept erasure under linear-affine predictors and squared loss, \citet{belrose2025leace} show that making a concept linearly unpredictable from a representation is equivalent to enforcing zero covariance between the representation and the concept label. \citet{belrose2025leace} formulated the following notion of guardedness:
% This gives a principled covariance-based view of concept control: concept erasure corresponds to removing the relevant cross-covariance.
% Thus, guardedness provides a principled justification for the covariance-removal constraint used in LEACE; the rest of the paper can then be read directly as extending this covariance-control view from erasure to switching.

\begin{definition}[Guardedness, \citealt{belrose2025leace}]
\label{def:guardedness}
Consider a $k$-class classification task over jointly defined random vectors $X$ and $Z$, where $X$ takes values in $\mathbb R^d$ and $Z$ is a one-hot label vector taking values in
\[
    \mathcal Z = \{\mathbf z \in \{0,1\}^k \mid \|\mathbf z\|_1 = 1\},
\]
with $\mathbb P(Z=j)>0$ for each class $j$. 
Let $\eta(\cdot;\vtheta):\mathbb R^d \to \mathbb R^k$ be a predictor from a function class
\[
    \mathcal V = \{\eta(\cdot;\vtheta)\mid \vtheta\in\Theta\},
\]
assumed to contain all constant functions, and let $\mathfrak L$ be a class of losses
$\loss:\mathbb R^k\times\mathcal Z\to[0,\infty)$. 
Let $\chi$ denote the set of random vectors of finite first moment taking values in $\mathbb R^d$ and jointly defined with $Z$.

We say that $X$ $(\mathcal V,\mathfrak L)$-\textbf{guards} $Z$ if, for every $\loss\in\mathfrak L$, it maximizes the minimum achievable prediction loss:
\begin{equation*}
    X \in 
    \mathop{\mathrm{argmax}}_{X'\in\chi}
    \inf_{\vtheta\in\Theta}
    \mathbb E\!\left[
        \loss\bigl(\eta(X';\vtheta), Z\bigr)
    \right].
\end{equation*}
Equivalently, $X$ is maximally uninformative about $Z$ for predictors in $\mathcal V$ and losses in $\mathfrak L$.
\end{definition}

The full guardedness definition is general, but the key consequence for our purposes is its linear characterization. 
For linear--affine predictors with squared loss, guardedness is equivalent to zero cross-covariance between the representation and the label.

\begin{theorem}[Linear guardedness, \citealt{belrose2025leace}]
\label{thm:linear_guardedness}
The following statements are equivalent:
\begin{itemize}
    \item $X$ linearly guards $Z$ in the sense of Definition~\ref{def:guardedness};
    \item every component of $X$ has zero covariance with every component of $Z$, i.e.
    \[
        \cov(X,Z)=0.
    \]
\end{itemize}
\end{theorem}

Therefore, in the affine setting, concept erasure can be expressed as a covariance-control problem: find a transformation of the representation that removes linear dependence on the concept label. 
At the same time, to preserve unrelated information, the transformation should perturb the representation as little as possible. 
These two desiderata lead to the LEACE objective: among all affine maps that enforce $\cov(AX+b,Z)=0$, choose the one closest to the identity transformation in expected squared distance.

\begin{theorem}[LEACE, \citealt{belrose2025leace}]
\label{thm:belrose}
Let $X$ and $Z$ be random vectors taking values in $\mathbb R^d$ and $\mathbb R^k$, respectively, each with finite second moments. 
Define $\Sigma_{XX}=\cov(X,X)\in\mathbb R^{d\times d}$, $\Sigma_{XZ}=\cov(X,Z)\in\mathbb R^{d\times k}$
and assume $\mathrm{Im}(\Sigma_{XZ})\subseteq\mathrm{Im}(\Sigma_{XX})$. 
Then the optimization problem
\begin{equation}
\label{eq:belrose_minimization}
    \min_{\substack{A\in\mathbb R^{d\times d}\\ b\in\mathbb R^d}}
    \mathbb E\!\left[
        \left\|AX+b-X\right\|_2^2
    \right]
    \quad
    \text{s.t.}
    \quad
    \cov(AX+b,Z)=0
\end{equation}
has the following solution, almost surely:
\begin{align}
    \widehat A
    &=
    I - W^+(W\Sigma_{XZ})(W\Sigma_{XZ})^+W,
    \label{eq:belrose_optimum_A}
    \\
    \widehat b
    &=
    \mathbb E[X] - \widehat A\,\mathbb E[X],
    \label{eq:belrose_optimum_b}
\end{align}
where $W=(\Sigma_{XX}^{1/2})^+$ is the whitening transformation.
\end{theorem}

Here and throughout, $A^+$ denotes the Moore--Penrose pseudo-inverse. 
For a positive semidefinite symmetric matrix $A=VSV^\top$, we write
$A^{1/2}=VS^{1/2}V^\top$, where $S^{1/2}$ is obtained by taking square roots of the diagonal entries of $S$.